% !TEX TS-program = lualatex

\documentclass[border = 3pt]{standalone}

\usepackage{main}

%%%
% Nous allons utiliser le "joli" vert ''ForestGreen''.
%%%
\usepackage[svgnames]{xcolor}

\begin{document}

%%%
% L'argument obligatoire donne le nombre de lignes attendu.
%
%
% note::
%     Pour obtenir une ligne monochrome, il faut forcément
%     dessiner `2^n` lignes.
%%%
\discretsierpinsky{16}

\quad

%%%
% Il est très pratique de disposer ici de deux types d'options
% n'utilisant pas les mêmes délimiteurs.
%
%     + La couleur des impairs est donnée entre crochets.
%
%     + La couleur des pairs est donnée entre chevrons.
%%%
\discretsierpinsky[orange]<ForestGreen>{32}

\end{document}
